% ============================================================
%  template.tex — Notas Acadêmicas
%  Autor: <Gabriel de Souza Dourado>
% ============================================================

\documentclass[12pt, a4paper]{article}

% ── Codificação e Língua ────────────────────────────────────
\usepackage[utf8]{inputenc}
\usepackage[T1]{fontenc}
\usepackage[brazil]{babel}

% ── Tipografia ──────────────────────────────────────────────
\usepackage{lmodern}
\usepackage{microtype}

% ── Margens ─────────────────────────────────────────────────
\usepackage[
  top=2.5cm, bottom=2.5cm,
  left=2.8cm, right=2.8cm
]{geometry}
\setlength{\skip\footins}{1.5em}

% ── Matemática ──────────────────────────────────────────────
\usepackage{amsmath, amssymb, amsthm}
\usepackage{mathtools}

% ── Algoritmos e Código ─────────────────────────────────────
\usepackage{algorithm}
\usepackage{algpseudocode}
\usepackage{listings}
\usepackage{xcolor}

% ── Pseudocódigo ─────────────────────────────────────
\renewcommand{\algorithmicrequire}{\textbf{Entrada:}}
\renewcommand{\algorithmicensure}{\textbf{Saída:}}

% ── Figuras e Tabelas ───────────────────────────────────────
\usepackage{graphicx}
\usepackage{booktabs}
\usepackage{caption}
\usepackage{float}

% ── Hiperlinks ──────────────────────────────────────────────
\usepackage[
  colorlinks=true,
  linkcolor=blue!70!black,
  citecolor=blue!70!black,
  urlcolor=blue!70!black
]{hyperref}

% ── Cabeçalho e Rodapé ──────────────────────────────────────
\usepackage{fancyhdr}
\pagestyle{fancy}
\fancyhf{}
\fancyhead[L]{\small\textit{\docsubtitle}}
\fancyhead[R]{\small\textit{\docauthor}}
\fancyfoot[C]{\small\thepage}
\renewcommand{\headrulewidth}{0.4pt}

% ── Cores ───────────────────────────────────────────────────
\definecolor{mainblue}{RGB}{30, 80, 160}
\definecolor{codegray}{RGB}{245, 245, 245}
\definecolor{codecomment}{RGB}{100, 100, 100}
\definecolor{codestring}{RGB}{163, 21, 21}
\definecolor{codekeyword}{RGB}{30, 80, 160}

% ── Estilo de Código (C++) ───────────────────────────────────
\lstdefinestyle{cppstyle}{
  language=C++,
  backgroundcolor=\color{codegray},
  basicstyle=\ttfamily\small,
  keywordstyle=\color{codekeyword}\bfseries,
  commentstyle=\color{codecomment}\itshape,
  stringstyle=\color{codestring},
  numberstyle=\tiny\color{codecomment},
  numbers=left,
  numbersep=8pt,
  breaklines=true,
  frame=single,
  framerule=0pt,
  rulecolor=\color{codegray},
  showstringspaces=false,
  tabsize=2,
  captionpos=b,
}
\lstset{style=cppstyle}

% ── Ambientes de Teorema ─────────────────────────────────────
\theoremstyle{definition}
\newtheorem{definition}{Definição}[section]
\newtheorem{example}{Exemplo}[section]

\theoremstyle{plain}
\newtheorem{theorem}{Teorema}[section]
\newtheorem{lemma}[theorem]{Lema}
\newtheorem{corollary}[theorem]{Corolário}
\newtheorem{proposition}[theorem]{Proposição}

\theoremstyle{remark}
\newtheorem*{remark}{Observação}
\newtheorem*{note}{Nota}

% ── Espaçamento ─────────────────────────────────────────────
\usepackage{setspace}
\setstretch{1.3}

% ============================================================
%  METADADOS — altere apenas estas 4 linhas em cada documento
% ============================================================
\newcommand{\docauthor}{Gabriel de Souza Dourado}
\newcommand{\doctitle}{Teoria da Computação}
\newcommand{\docsubtitle}{Autômatos Finitos Não Determinísticos}
\newcommand{\docdate}{\today}
% ============================================================

\begin{document}

% ── Título ───────────────────────────────────────────────────
\begin{center}
  {\large \textcolor{mainblue}{\textsc{\docsubtitle}}}\\[0.4em]
  {\LARGE \textbf{\doctitle}}\\[0.6em]
  {\small \docauthor \quad·\quad \docdate}
  \vspace{0.4em}
  \hrule
\end{center}

\vspace{1em}

% ── Sumário (descomente se quiser) ───────────────────────────
% \tableofcontents
% \newpage

% ------------------------------------------------------------
\section{Introdução}
% ------------------------------------------------------------

Um \textbf{Autômato Finito Não Determinístico (AFND)}, assim como um AFD, é um modelo matemático de computação utilizado para reconhecer linguagens regulares, porém, não possui do determinismo na suya função de transição. Formalmente, um AFND é caracterizado por cinco componentes, assim como um AFD:

\begin{enumerate}
    \item Um conjunto finito de \textbf{estados} denotado $Q$.
    \item Um conjunto finito de \textbf{símbolos de entrada} denotado por $\Sigma$.
    \item Uma \textbf{função de transição} $\delta$\footnote{É na função de transição que está a diferença entre um AFD e um AFND.}.
    \item Um \textbf{estado inicial} $q_0$, estado onde o autômato começa a computação.
    \item Um subconjunto de $Q$, denotado de $F$, de \textbf{estados de aceitação}.
\end{enumerate}

\section{Definição}

\begin{definition}[Autômato Finito Não Determinístico]
É dado pela 5-tupla ordenada:

\[
M = (Q, \Sigma, \delta, q_0, F),
\]
onde:

\begin{itemize}
    \item $Q$ é um conjunto finito de estados;
    \item $\Sigma$ é o alfabeto de entrada;
    \item $\delta : Q \times \Sigma \to 2^{Q}$ é a função de transição total;
    \item $q_0 \in Q$ é o estado inicial;
    \item $F \subseteq Q$ é o conjunto de estados de aceitação.
\end{itemize}

\end{definition}

% ------------------------------------------------------------
\section{Implementação de um AFND (Backtracking)}
% ------------------------------------------------------------
O não-determinismo da função de transição $\delta$ aumenta consideravelmente a complexidade de implementação em relação ao AFD. Aqui, adotamos a estratégia de \textit{backtracking}: o algoritmo avança recursivamente pelos caminhos de computação do AFND e, ao encontrar uma configuração sem transição definida ou um estado não-aceitador ao fim da palavra, retrocede e explora as demais transições disponíveis. A palavra é aceita se \textbf{algum} caminho levar a um estado de aceitação após consumir toda a entrada.

\begin{algorithm}[H]
  \caption{$\textsc{Busca}(M, w, q, i)$}
  \label{alg:afndBusca}
  \begin{algorithmic}[1]
    \Require Um AFND $M = (Q, \Sigma, \delta, q_0, F)$,
             onde $\delta : Q \times \Sigma \to 2^Q$,
             uma palavra $w \in \Sigma^*$,
             um estado $q \in Q$ e um índice $i \in \mathbb{N}$
    \Ensure  \textbf{verdadeiro} se existe um caminho de aceitação, 
             \textbf{falso} caso contrário
    \If{$i = |w|$}
      \If{$q \in F$}
        \State \Return \textbf{verdadeiro}
      \Else
        \State \Return \textbf{falso}
      \EndIf
    \EndIf
    \For{cada $q' \in \delta(q,\ w[i])$}
      \If{$\textsc{Busca}(M, w, q', i + 1)$}
        \State \Return \textbf{verdadeiro}
      \EndIf
    \EndFor
    \State \Return \textbf{falso}
  \end{algorithmic}
\end{algorithm}

\begin{algorithm}[H]
  \caption{Simulação de um AFND com Backtracking}
  \label{alg:afndBacktracking}
  \begin{algorithmic}[1]
    \Require Um AFND $M = (Q, \Sigma, \delta, q_0, F)$,
            onde $\delta : Q \times \Sigma \to 2^Q$,
             e $w \in \Sigma^*$
    \Ensure  \textbf{aceita} se $w \in L(M)$, \textbf{rejeita} caso contrário
    \If{$\textsc{Busca}(M, w, q_0, 0)$}
      \State \Return \textbf{aceita}
    \Else
      \State \Return \textbf{rejeita}
    \EndIf
  \end{algorithmic}
\end{algorithm}

% ============================================================
%  REFERÊNCIAS
% ============================================================
\newpage
\begin{thebibliography}{9}

\bibitem{menezes}
MENEZES, P. B.\@.
\textit{Linguagens Formais e Autômatos}.
6.~ed. Porto Alegre: Bookman, 2011.

\bibitem{hopcroft}
HOPCROFT, J. E.; MOTWANI, R.; ULLMAN, J. D.\@.
\textit{Introduction to Automata Theory, Languages, and Computation}.
3.~ed. Boston: Pearson, 2006.

\end{thebibliography}

\end{document}